\section{Introduction}
\label{sec:introduction}

Service Network Design (SND) is the problem of determining service plans---routes, schedules, and resource capacities---to satisfy transportation demand over a network. It occupies a central position in transportation planning and operations research, as the service plan dictates both the operating cost structure and the level of service experienced by customers. Foundational SND models cast the problem as a multi-commodity flow (MCF) optimization on time-space networks, where nodes represent locations at specific points in time and arcs represent either service operations or idle resource holding \citep{boland2017continuous}. The resulting mixed-integer programs (MIPs) must simultaneously determine which services to operate, when to operate them, and which resource types to assign, making SND one of the most challenging combinatorial optimization problems in practice.

A critical limitation of classical SND formulations is the assumption that demand is exogenous---fixed independently of the service plan. In reality, when travelers can choose among competing itineraries offered by one or more operators, or opt out of traveling entirely, the demand becomes endogenous: it depends on the service plan itself. Incorporating passenger choice models into SND creates a feedback loop between supply and demand decisions. The service plan determines the set of available itineraries, which in turn determines the demand captured, which in turn affects the profitability of the service plan. This coupling makes the problem significantly harder---the feasible region now includes choice-consistency constraints that link operational and demand variables---but also far more realistic and valuable for practical planning.

The integration of SND with passenger choice arises naturally in several transportation domains. In \textit{airline scheduling}, carriers jointly optimize fleet assignment, timetabling, and route planning while accounting for how passengers select among available itineraries based on schedule convenience, connection quality, and fare levels \citep{wei2020airline, yan2022choice}. In \textit{urban rail and metro operations}, transit agencies design train frequencies and stopping patterns on networks where commuters choose routes based on travel time, number of transfers, and crowding levels. In \textit{intercity bus and rail services}, operators plan service frequencies and departure schedules for corridors where travelers choose between competing operators and modes based on departure times, travel duration, and ticket price. Across all these settings, the core modeling challenge is the same: the combinatorial explosion of feasible itineraries over the time-space network, coupled with the need to endogenize demand through a discrete choice framework.

The standard modeling approach for SND with passenger choice formulates the problem as a multi-commodity flow program on a time-space network with an embedded discrete choice model. The planning horizon is discretized into time intervals, and service arcs are defined over this discretized network. Passenger demand is endogenized through the Sales-Based Linear Program (SBLP) linearization of the General Attraction Model (GAM) \citep{gallego2015general}, which provides a tractable linear representation of choice probabilities within the MIP. The resulting formulation simultaneously optimizes service operations (binary assignment of resources to arcs) and demand allocation (continuous market share variables governed by choice-consistency constraints). However, the model is extremely large: even for moderately sized networks, the number of feasible itineraries can reach hundreds of thousands, and the interaction between binary service variables and continuous choice variables produces a dense, ill-conditioned constraint matrix.

Existing solution methods face fundamental limitations when applied to this problem class. Dynamic Discretization Discovery (DDD), introduced by \citet{boland2017continuous}, efficiently compresses the time-space network by starting with a coarse discretization and iteratively refining only where needed, guaranteeing convergence to the continuous-time optimum. However, DDD was designed for freight and service network problems with exogenous demand; it lacks the mechanism to handle the itinerary-based passenger choice behavior that couples demand allocation to the service plan. Decomposition methods have been applied to improve the tractability of choice-based airline schedule-design models \citep{yan2022choice}. However, these approaches are primarily designed to produce high-quality heuristic solutions and do not provide an exact dynamic-refinement framework for the itinerary-level, resource-coupled formulation studied here. Direct MIP solving of the full formulation is intractable for realistic instances: the massive number of itinerary variables---potentially hundreds of thousands to millions---overwhelms commercial solvers, which stall in the branch-and-bound process due to weak LP relaxations and fractional degeneracy.

In this paper, we propose an exact solution framework that overcomes these limitations. Our main contributions are as follows:
\begin{enumerate}
    \item We develop an \textit{Itinerary Aggregation} technique that, combined with time discretization, produces a tight lower bound by constructing ``super-itineraries.'' These aggregated itineraries eliminate the fictitious revenue that arises from fractional resource allocations in the relaxed network, yielding a relaxation that is substantially tighter than the standard LP relaxation of the full model.
    \item We design a \textit{repair algorithm} for upper bound generation, based on frequency fixing and itinerary cuts, that reconstructs feasible integer solutions from the relaxed lower-bound solution. This structure-aware approach exploits the information embedded in the relaxation to produce high-quality feasible schedules.
    \item We propose a \textit{refinement strategy} that dynamically adds time points to the discretized network by detecting and resolving three types of relaxation artifacts---chronological violations, separation infeasibilities, and phantom revenue---guaranteeing finite convergence to the global optimum.
    \item We validate the framework on real-world instances, demonstrating that previously intractable problems with nearly half a million itinerary variables can be solved to provable optimality.
\end{enumerate}

The remainder of this paper is organized as follows. Section~\ref{sec:literature_review} reviews the relevant literature on service network design and dynamic discretization discovery. Section~\ref{sec:model_overview} formally describes the problem and presents the mathematical formulation on the full time-space network. Section~\ref{sec:algorithm} details the solution methodology, including the relaxation model that provides a valid lower bound (Section~\ref{sec:lbm}), the upper bound generation procedure (Section~\ref{sec:ub}), and the refinement strategy that guarantees convergence (Section~\ref{sec:refinement}). Section~\ref{sec:comp_study} presents the computational study, demonstrating the practical applicability of the framework. We conclude with a summary of findings and directions for future research.
